% Generated draft: requires mathematical and source review.
\documentclass[11pt]{article}
\usepackage[margin=25mm]{geometry}
\usepackage{fontspec}
\setmainfont{Latin Modern Roman}
\usepackage{amsmath,amssymb,mathtools,unicode-math}
\setmathfont{Latin Modern Math}
\usepackage{enumitem,needspace,xurl,hyperref}
\setlength{\parindent}{0pt}
\setlength{\parskip}{5pt}
\setlength{\emergencystretch}{3em}
\newcommand{\N}{\mathbb{N}}
\newcommand{\Z}{\mathbb{Z}}
\newcommand{\Q}{\mathbb{Q}}
\newcommand{\R}{\mathbb{R}}
\newcommand{\C}{\mathbb{C}}
\newcommand{\F}{\mathbb{F}}
\newcommand{\A}{\mathbb{A}}
\newcommand{\PP}{\mathbb{P}}
\newcommand{\E}{\mathbb{E}}
\newcommand{\eps}{\varepsilon}
\newcommand{\dd}{\,\mathrm{d}}
\begin{document}
\section*{UnsolvedMath ID: 30004797}
\textbf{Source problem code:} catalogue entry
\subsection{Definitions and mathematical statement}
The supplied record does not specify enough data to state a uniquely determined Plancherel-measure identity. The following gives a conditional formulation of its intended target and identifies the missing qualifications.

In the usual nonarchimedean representation-theoretic setting, let \(F\) be a nonarchimedean local field, let \(G\) be a connected reductive group over \(F\), and let \(\widehat{G(F)}_{\mathrm{temp}}\) denote the set of equivalence classes of irreducible tempered representations of \(G(F)\). For a parabolic subgroup \(P=MN\), with Levi factor \(M\) and unipotent radical \(N\), write
\[
I_P^G(\sigma)=\operatorname{Ind}_{P(F)}^{G(F)}\bigl(\delta_P^{1/2}\sigma\bigr)
\]
for normalized parabolic induction, where \(\delta_P\) is the modulus character and \(\sigma\) is inflated from \(M(F)\) to \(P(F)\). This specifies notation in a conventional setting; the record does not establish the intended range of fields and groups or the additional hypotheses required for Yu's construction.

\emph{Requested assertion, subject to the unresolved specifications below.} For every nondiscrete component of \(\widehat{G(F)}_{\mathrm{temp}}\) arising by parabolic induction from nonsingular Yu supercuspidal representations of Levi subgroups, verify the Plancherel-measure formula referred to in the source as the Hiraga--Ichino--Ikeda Plancherel-measure conjecture.

Here the terms ``nonsingular Yu supercuspidal'' and ``nondiscrete component'' retain their source-specific meanings. The supplied record gives neither their precise conventions nor the conjectured density formula. In particular, it supplies no parameterization of the relevant components, Haar-measure normalization, reference measure for a density, or complete hypotheses. Consequently, the requested assertion is a qualified target for formalization, rather than a fully specified measure identity; these omissions must be resolved from the cited source before a precise mathematical statement can be finalized.
\subsection{Short English statement}
The problem asks for the Hiraga--Ichino--Ikeda Plancherel-measure formula on the nondiscrete tempered components obtained by inducing nonsingular Yu supercuspidal representations. The supplied record omits the formula and essential hypotheses and normalizations, so a complete precise conjecture cannot yet be extracted. Its open status is supplied dated curation, not an independently verified conclusion.
\subsection{Sources}
\begin{itemize}
\item Local registry: \texttt{data/unsolvedmath/problems.json}, permanent UnsolvedMath ID \(30004797\), source problem code \texttt{catalogue entry}, title ``Plancherel Densities on Yu-Induced Tempered Components,'' in the supplied Oberwolfach Reports Open Problems collection.
\item Supplied primary-source citation: \emph{Automorphic Forms, Geometry and Arithmetic (2022)}, Oberwolfach Reports, DOI \texttt{10.4171/owr/2021/39}, pp.\ 2089--2156. \url{https://doi.org/10.4171/owr/2021/39}. The record assigns proposed year 2021 and supplies no individual proposer or exact problem location within the report.
\item Supplied literature source, labeled ``formal degrees and parabolic induction'': \url{https://doi.org/10.1007/s00229-025-01632-z}.
\item Supplied literature source, labeled ``affine Hecke algebras and HII Plancherel conjectures'': \url{https://arxiv.org/abs/1807.10232}.
\item Literature assessment dated exactly \texttt{2026-08-22}: the supplied curation labels the target open and states that recent work addresses only narrower generic or otherwise special parabolic-induction cases. This is supplied dated curation derived from the OWR record, not independently verified literature assessment; the time-sensitive status and cited claims require independent verification.
\item The supplied \texttt{corrected\_verified} statement status is explained as a publication-format audit that preserved mathematical content while repairing notation or typography. It does not establish mathematical correctness or independent source verification.
\end{itemize}
\end{document}
